\section[Cote 88 : l'ordre $2\nu$ et la trichotomie des cartes standard]{Cote 88 : l'ordre $2\nu$ et la trichotomie des cartes standard\protect\verifie{Folder88}}\label{sec:88}

La cote 88, \emph{Cartes standards : notes manuscrites (s.d.)}, quinze pages sans date, appartient au groupe \q{Géométrie et topologie combinatoire} que l'inventaire date \q{1976-[vers 1986]}. Elle contient deux textes. Le premier, sous le titre de chemise \q{Cartes associées aux syst.\ de Racines}, remplace un système de racines par son groupe de Weyl, le torseur de ses bases et ses réflexions simples, fait opérer des involutions $\sigma_0, \dots, \sigma_{\ell-1}$ sur les drapeaux du complexe qu'il en tire, et calcule à la page~8 l'ordre du dernier produit : \q{ordre de $\pi = \sigma_{\ell-2}\,\sigma_{\ell-1} = 2\nu$}, $\nu$ étant \q{le ppcm des $n_{ij}$} (\lot{88}{1}{8}). Le second, intitulé de sa main \q{Théorème (à prouver)}, énonce en sept articles la théorie des cartes régulières de type $(p, q)$ ; son article~(6), à la page~14, les répartit en sphériques, euclidiennes et hyperboliques par des listes de couples (\lot{88}{1}{14}). Deux énoncés de la lecture sont démontrés ici : la proposition de la page~8, et l'arithmétique qui trie les cas de l'article~(6).

\begin{definition}
Soient $W$ un groupe, $R$ un ensemble non vide sur lequel $W$ opère librement et transitivement, $I$ un ensemble fini de cardinal $\ell \geq 2$, et $\varphi : R \times I \to W$ une application telle que
\[
  \text{a)}\ \ \varphi(wr, i) = w\,\varphi(r, i)\,w^{-1}, \qquad \varphi(r, i)^2 = 1 \qquad (w \in W,\ r \in R,\ i \in I).
\]
Soit $\mathrm{Rep}'(I)$ l'ensemble des bijections $u : [1, \ell] \to I$, et $\tau$ la transposition de $\ell - 1$ et $\ell$. On définit sur $R \times \mathrm{Rep}'(I)$
\[
  \sigma_{\ell-2}(r, u) = (r,\ u \circ \tau), \qquad
  \sigma_{\ell-1}(r, u) = \bigl(\varphi(r, u(\ell))\, r,\ u\bigr).
\]
Pour $g \in W$, on note $\operatorname{ord}(g)$ le générateur positif ou nul du sous-groupe $\{k \in \mathbf Z \mid g^k = 1\}$, qui est nul si $g$ est d'ordre infini, et de même pour une permutation. On pose $n_{ij} = \operatorname{ord}\bigl(\varphi(r, i)\varphi(r, j)\bigr)$ et $\nu = \operatorname{ppcm}\{n_{ij} \mid i \neq j\}$, un ppcm dont un terme est nul étant nul.
\end{definition}

\noindent Dans la lecture, $W$ est le groupe de Weyl, $R$ l'ensemble des bases, $\varphi(r, i)$ la réflexion simple d'indice $i$ relative à la base $r$ ; seules sont retenues ici, des quatre conditions a) à d) que la lecture énonce, la condition a) et la moitié $\varphi(r,i)^2 = 1$ de c).

\begin{enonce}[Proposition, page 8]
L'ordre de $\sigma_{\ell-2}\sigma_{\ell-1}$ sur $R \times \mathrm{Rep}'(I)$ est $2\nu$, où $\nu$ est le plus petit commun multiple des $n_{ij}$, $i \neq j$.
\end{enonce}

\begin{enonce}[Article (6), page 14 : l'arithmétique]
Soient $p, q \geq 1$ des entiers, le type $(p, q)$ étant atteint au sens de l'article~(2) : on exclut $(1, q)$ pour $q \neq 2$ et $(p, 1)$ pour $p \neq 2$. Alors
\begin{enumerate}
\item $\frac1p + \frac1q > \frac12$ si et seulement si $p = 2$, ou $q = 2$, ou $(p, q)$ est l'un de $(3,3)$, $(3,4)$, $(4,3)$, $(3,5)$, $(5,3)$ ;
\item $\frac1p + \frac1q = \frac12$ si et seulement si $(p, q)$ est l'un de $(4,4)$, $(3,6)$, $(6,3)$ ;
\item pour $p, q \geq 3$, $\frac1p + \frac1q < \frac12$ si et seulement si $p + q \geq 9$ et $(p, q) \neq (3,6), (6,3)$.
\end{enumerate}
\end{enonce}

On pose $\pi = \sigma_{\ell-1}\sigma_{\ell-2}$, le produit avec lequel calcule la page~6, et, pour $(r, u)$ donné, $a = \varphi(r, u(\ell-1))$ et $b = \varphi(r, u(\ell))$.

\begin{lemme}\label{lem:88-invol}
$\sigma_{\ell-2}$ et $\sigma_{\ell-1}$ sont des involutions ; par suite $\sigma_{\ell-2}\sigma_{\ell-1} = \pi^{-1}$.
\end{lemme}

\begin{proof}
Pour $\sigma_{\ell-2}$, c'est $\tau^2 = \mathrm{id}$. Pour $\sigma_{\ell-1}$, posons $s = \varphi(r, i)$ avec $i = u(\ell)$ ; par a), $\varphi(sr, i) = s\,s\,s^{-1} = s$, de sorte que $\sigma_{\ell-1}^2(r, u) = (s^2 r, u) = (r, u)$. Alors $\sigma_{\ell-2}\sigma_{\ell-1} = \sigma_{\ell-2}^{-1}\sigma_{\ell-1}^{-1} = (\sigma_{\ell-1}\sigma_{\ell-2})^{-1}$.
\end{proof}

\begin{lemme}\label{lem:88-carre}
On a $\pi(r, u) = \bigl(a\, r,\ u \circ \tau\bigr)$ ; pour tout $g \in W$, $\pi^2(g r, u) = (g\,ab\,r, u)$ ; et, pour tout entier $m \geq 0$, $\pi^{2m}(r, u) = \bigl((ab)^m r,\ u\bigr)$ et la seconde composante de $\pi^{2m+1}(r, u)$ est différente de $u$.
\end{lemme}

\begin{proof}
Comme $(u \circ \tau)(\ell) = u(\ell - 1)$, $\pi(r, u) = \sigma_{\ell-1}(r, u \circ \tau) = (a r, u \circ \tau)$. Pour $g \in W$, a) donne $\varphi(gr, u(\ell-1)) = g a g^{-1}$, d'où $\pi(gr, u) = (g a r, u \circ \tau)$ ; puis, $(u \circ \tau)(\ell - 1) = u(\ell)$ et $(u \circ \tau) \circ \tau = u$, et a) donne $\varphi(gar, u(\ell)) = (ga)\,b\,(ga)^{-1}$, d'où $\pi(gar, u \circ \tau) = (g a b r, u)$. La formule pour $\pi^{2m}$ s'en déduit par récurrence, en appliquant la précédente à $g = (ab)^m$. Enfin $\pi^{2m+1}(r, u) = \pi\bigl((ab)^m r, u\bigr)$ a pour seconde composante $u \circ \tau$, et $(u \circ \tau)(\ell) = u(\ell - 1) \neq u(\ell)$ puisque $u$ est injective.
\end{proof}

\begin{lemme}\label{lem:88-fixe}
Pour tout entier $k \geq 0$, $\pi^k(r, u) = (r, u)$ si et seulement si $2\, n_{u(\ell-1),\, u(\ell)}$ divise $k$.
\end{lemme}

\begin{proof}
Si $k = 2m$ est pair, $\pi^k(r, u) = ((ab)^m r, u)$ par le lemme~\ref{lem:88-carre} ; l'action étant libre, c'est $(r, u)$ si et seulement si $(ab)^m = 1$, c'est-à-dire si et seulement si $\operatorname{ord}(ab)$ divise $m$, soit $2\operatorname{ord}(ab)$ divise $k$. Si $k$ est impair, $\pi^k(r, u) \neq (r, u)$ par le même lemme, et $2\operatorname{ord}(ab)$ ne divise pas $k$. Enfin $ab = \varphi(r, u(\ell-1))\varphi(r, u(\ell))$ est d'ordre $n_{u(\ell-1), u(\ell)}$ par le lemme suivant, qui assure que cet ordre ne dépend pas de $r$.
\end{proof}

\begin{lemme}\label{lem:88-indep}
L'ordre de $\varphi(r, i)\varphi(r, j)$ ne dépend pas de $r \in R$. Et pour $i \neq j$ dans $I$, il existe $u \in \mathrm{Rep}'(I)$ avec $u(\ell - 1) = i$ et $u(\ell) = j$.
\end{lemme}

\begin{proof}
Soient $r, r' \in R$ ; par transitivité, $r' = wr$ pour un $w \in W$, et par a), $\varphi(r', i)\varphi(r', j) = w\,\varphi(r, i)\varphi(r, j)\,w^{-1}$, qui est conjugué de $\varphi(r, i)\varphi(r, j)$ et a donc le même ordre. Pour la seconde assertion, partons d'une bijection $e : [1, \ell] \to I$, qui existe puisque $I$ a $\ell$ éléments. On la compose à gauche par la transposition de $e(\ell - 1)$ et de $i$, pour obtenir $u_1$ avec $u_1(\ell - 1) = i$, puis par la transposition de $u_1(\ell)$ et de $j$. Comme $u_1(\ell) \neq i$ et $j \neq i$, cette seconde transposition fixe $i$, et la bijection obtenue envoie $\ell - 1$ sur $i$ et $\ell$ sur $j$.
\end{proof}

\begin{proof}[Démonstration de la proposition]
Par le lemme~\ref{lem:88-invol}, $\sigma_{\ell-2}\sigma_{\ell-1}$ a le même ordre que $\pi$. Montrons que $\pi^k = \mathrm{id}$ si et seulement si $2n_{ij}$ divise $k$ pour tous $i \neq j$. Si $\pi^k = \mathrm{id}$ et $i \neq j$, prenons $u$ comme au lemme~\ref{lem:88-indep} et un $r$ quelconque de $R$, qui n'est pas vide : le lemme~\ref{lem:88-fixe} en $(r, u)$ donne $2n_{ij} \mid k$. Réciproquement, pour tout $(r, u)$, les indices $u(\ell - 1)$ et $u(\ell)$ sont distincts, et le lemme~\ref{lem:88-fixe} donne $\pi^k(r, u) = (r, u)$. Or $2n_{ij} \mid k$ pour tous $i \neq j$ équivaut à $2\nu \mid k$ : si c'est le cas pour tous, $k$ est pair, puisqu'il existe au moins un couple $i \neq j$, et $k/2$ est un multiple commun des $n_{ij}$, donc de $\nu$ ; la réciproque vient de ce que chaque $n_{ij}$ divise $\nu$. Ainsi $\{k \mid \pi^k = \mathrm{id}\} = 2\nu\,\mathbf Z$, et l'ordre de $\pi$ est $2\nu$.
\end{proof}

\begin{proof}[Démonstration de l'arithmétique de (6)]
Pour $p, q \geq 1$, $\frac1p + \frac1q > \frac12$, $= \frac12$, $< \frac12$ équivalent respectivement à $pq < 2(p + q)$, $pq = 2(p + q)$, $pq > 2(p + q)$, c'est-à-dire à $(p-2)(q-2) < 4$, $= 4$, $> 4$. Si $p \geq 7$, $pq < 2(p+q)$ entraîne $q \leq 2$, et symétriquement ; hors de ces cas, il reste un nombre fini de couples avec $p, q \leq 6$, qu'on examine un à un. Pour~(1), $p = 2$ ou $q = 2$ donnent $0 < 4$ ; $p = 1$ n'est permis qu'avec $q = 2$, et $q = 1$ qu'avec $p = 2$ ; pour $p, q \geq 3$, $(p-2)(q-2) \in \{1, 2, 3\}$ laisse exactement $(3,3)$, $(3,4)$, $(4,3)$, $(3,5)$, $(5,3)$. Pour~(2), $p$ ou $q$ égal à $1$ ou $2$ est exclu, puisque $(p-2)(q-2)$ vaut alors $0$, ou $2 - q$ ou $2 - p$, qui est au plus $1$ ; pour $p, q \geq 3$, $(p-2)(q-2) = 4$ donne $(3,6)$, $(4,4)$, $(6,3)$. Pour~(3), les couples $p, q \geq 3$ qui ne vérifient pas $(p-2)(q-2) > 4$ sont les huit couples obtenus en~(1) et~(2). Six d'entre eux, les cinq de~(1) et $(4,4)$, sont exactement les couples $p, q \geq 3$ de somme au plus~$8$ ; les deux autres, $(3,6)$ et $(6,3)$, sont de somme~$9$. Un couple $p, q \geq 3$ vérifie donc $(p-2)(q-2) > 4$ si et seulement si $p + q \geq 9$ et $(p, q) \neq (3,6), (6,3)$.
\end{proof}

\begin{remarque}
La proposition ne demande qu'un groupe $W$ opérant librement et transitivement sur $R$ et une application $\varphi$ vérifiant a) et $\varphi(r, i)^2 = 1$ : que $W$ soit un groupe de Weyl, la condition b), la présentation c) et le graphe de Coxeter d) n'interviennent pas.
\end{remarque}

\begin{remarque}
La finitude de $W$ n'intervient pas : avec la convention adoptée pour l'ordre, la formule $2\nu$ vaut encore quand un $n_{ij}$ est infini.
\end{remarque}

\begin{remarque}
La page~14 écrit la condition $(c_2)$ \q{$p + q \geqslant 10$ ou $(p,q) \neq (3,6), (6,3)$} ; la forme juste, que la lecture rétablit et qui est démontrée en~(3), est $p + q \geq 9$ \emph{et} $(p, q) \neq (3,6), (6,3)$.
\end{remarque}

\begin{remarque}
Dans la lecture, la présentation de Coxeter du groupe cartographique, $\widehat\Gamma = \langle \sigma_0, \sigma_1, \sigma_2 \mid \sigma_i^2,\ (\sigma_0\sigma_2)^2 \rangle$, est une définition reprise de l'\emph{Esquisse d'un programme}, non un énoncé.
\end{remarque}

\noindent Ne sont pas démontrés ici : la proposition des pages~3 et~4 ($\mathrm{Aut}(\Sigma) \to \mathrm{Aut}(\Phi(\Sigma))$ est injectif, et bijectif exactement hors des types $B_2$, $G_2$, $F_4$ et des couples $B_n$, $C_n$) ; les autres relations entre les $\sigma_i$ de la page~6 ; la triangulation de la sphère $S^{\ell-1}$ par le complexe des drapeaux ; les articles~(1) à~(5) et~(7) du second texte (régularité des cartes quasi-régulières simplement connexes, unicité de la carte standard de type $(p,q)$, $\mathrm{Aut}(C_{p,q}) \simeq \Gamma_{p,q}$, réalisation géodésique sur les trois surfaces standard, cartes de type $(p,q)$ comme quotients) ; le contenu géométrique de~(6), finitude de $\Gamma_{p,q}$ dans le cas sphérique et sous-groupes libres dans le cas hyperbolique.
